% ============================================================================
%  CAPÍTULO IV: MEDIDAS DE MITIGACIÓN Y ENDURECIMIENTO TÉCNICO
% ============================================================================

\section{Estrategia Integral de Remediación y Endurecimiento}

Para resolver de manera definitiva las vulnerabilidades identificadas en el capítulo anterior, se diseñó e implementó un plan de mitigación técnica basado en el principio de defensa en profundidad (\textit{Defense in Depth}) \citep{cis_debian_2023}. A continuación se detallan las medidas aplicadas en cada nivel de la infraestructura.

\section{Remediación por Vulnerabilidad Identificada}

\subsection{Solución VULN-01: Deshabilitación de Suites Anónimas Diffie-Hellman en Postfix}
Para erradicar la vulnerabilidad de Man-in-the-Middle en los puertos de correo \texttt{25} (SMTP) y \texttt{587} (Submission):

\begin{enumerate}
  \item \textbf{Exclusión de Cifrados Anónimos (\textit{aNULL}):} En el archivo de configuración del servidor de correo Postfix (\texttt{services/mail/config/main.cf}), se excluyeron explícitamente todas las suites de cifrado anónimas, obsoletas o sin autenticación:
\begin{lstlisting}[style=terminal, caption={Endurecimiento Criptografico TLS en Postfix (main.cf)}, label={lst:mitigacion-postfix}]
# Excluir suites anonimas (aNULL), nulas y algoritmos debiles
smtpd_tls_exclude_ciphers = aNULL, eNULL, EXPORT, DES, RC4, MD5, PSK, aECDH, EDH-DSS-DES-CBC3-SHA, EDH-RSA-DES-CBC3-SHA, KRB5-DES, kDHr, kDHd, kEDH

# Forzar cifrado de alta seguridad
smtpd_tls_ciphers = high
smtpd_tls_mandatory_ciphers = high
smtpd_tls_protocols = !SSLv2, !SSLv3, !TLSv1, !TLSv1.1
smtpd_tls_mandatory_protocols = !SSLv2, !SSLv3, !TLSv1, !TLSv1.1

# Generar y vincular parametros Diffie-Hellman seguros de 2048 bits
smtpd_tls_dh1024_param_file = /etc/postfix/dh2048.pem
\end{lstlisting}

  \item \textbf{Generación de Parámetros DH Fuertes:} Se generó un archivo de parámetros Diffie-Hellman de 2048 bits mediante OpenSSL para reemplazar los grupos débiles de 1024 bits:
\begin{lstlisting}[style=terminal, caption={Generacion de Parametros Diffie-Hellman de 2048 bits}, label={lst:dhparam}]
openssl dhparam -out services/mail/config/dh2048.pem 2048
\end{lstlisting}
\end{enumerate}

\subsection{Solución VULN-02: Aislamiento del Motor PostgreSQL en Capa de Red}
Para neutralizar la exposición directa de la base de datos (puerto 5432):

\begin{enumerate}
  \item \textbf{Aislamiento en Docker Bridge:} Se modificó el archivo \texttt{docker-compose.yml} para remover la publicación pública del puerto \texttt{5432:5432}, sustituyéndola por la cláusula \texttt{expose: - 5432}. Con esto, el motor PostgreSQL únicamente es accesible para los contenedores que coexisten en la red interna de Docker (\texttt{homelab\_net}).
  \item \textbf{Filtrado en Firewall Nftables:} Se configuró una regla restrictiva en \texttt{server/nftables.conf} que vincula el puerto 5432 exclusivamente al conjunto de IPs de administradores autorizados (\texttt{@admin\_ips}):
\begin{lstlisting}[style=terminal, caption={Regla de Filtrado para PostgreSQL en nftables.conf}, label={lst:mitigacion-postgres}]
# Permitir acceso al motor PostgreSQL UNICAMENTE desde estaciones admin
tcp dport 5432 ip saddr @admin_ips accept comment "Postgres solo admins"
tcp dport 5432 drop comment "Bloqueo por defecto de Postgres para la LAN"
\end{lstlisting}
\end{enumerate}

\subsection{Solución VULN-03: Desactivación de Modo Debug y Divulgación en Webmail}
Para resolver la advertencia de \texttt{DEBUG is enabled} en el puerto 8080:
\begin{itemize}
  \item En el archivo de configuración \texttt{services/webmail/config/config.inc.php}, se establecieron los parámetros de depuración en falso:
\begin{lstlisting}[style=terminal, caption={Desactivacion de Modo Debug en Roundcube}, label={lst:mitigacion-debug}]
$config['debug_level'] = 1;
$config['devel_mode'] = false;
$config['display_errors'] = false;
\end{lstlisting}
  \item En Nginx se configuró la directiva \texttt{server\_tokens off;} para suprimir la divulgación de banners de software en cabeceras HTTP.
\end{itemize}

\subsection{Solución VULN-04: Cifrado TLS Obligatorio y Terminación SSL en Nginx}
Para erradicar la transmisión de credenciales y cookies de sesión en texto claro:
\begin{enumerate}
  \item Se centralizó la terminación criptográfica en el proxy inverso Nginx, configurando certificados TLS modernos con protocolos TLSv1.2 y TLSv1.3 y deshabilitando suites obsoletas (RC4, 3DES, CBC).
  \item Se implementó redirección automática en el puerto 80 hacia HTTPS (\texttt{return 301 https://\$host\$request\_uri;}).
  \item Los servicios web internos (Roundcube en 8080, Filemanager en 8082 y Grafana en 3000) se enrutan de forma segura bajo canales cifrados.
\end{enumerate}

\subsection{Solución VULN-05: Endurecimiento de Cabeceras Web y WAF Anti-SQLi}
En el archivo \texttt{services/web/nginx.conf} se añadieron las siguientes directivas de endurecimiento bancario y reglas de inspección profunda de parámetros:

\begin{lstlisting}[style=terminal, caption={Cabeceras de Endurecimiento Web y Bloqueo Anti-SQLi en Nginx}, label={lst:mitigacion-headers}]
# Cabeceras de seguridad bancaria OWASP / PCI-DSS
add_header X-Frame-Options "DENY" always;
add_header X-Content-Type-Options "nosniff" always;
add_header X-XSS-Protection "1; mode=block" always;
add_header Strict-Transport-Security "max-age=31536000; includeSubDomains" always;

# Bloqueo proactivo en tiempo real contra Inyecciones SQL
if ($bloqueo_sqli) {
    return 403 "Acceso denegado [WAF-ANTI-SQLI]: Carga util de Inyeccion SQL detectada y neutralizada por el Firewall Web.\n";
}

# Bloqueo proactivo contra Path Traversal / LFI
if ($bloqueo_traversal) {
    return 403 "Acceso denegado [WAF-ANTI-LFI]: Intento de Directory Traversal detectado y neutralizado.\n";
}
\end{lstlisting}

\subsection{Solución VULN-06: Mitigación Activa contra DoS, Slowloris y Tarpitting Defensivo}
Para prevenir la saturación de recursos y neutralizar escaneos masivos:

\begin{enumerate}
  \item \textbf{Timeouts y Rate Limiting en Nginx:}
\begin{lstlisting}[style=terminal, caption={Directivas Anti-Slowloris y Limitacion de Conexiones en Nginx}, label={lst:mitigacion-slowloris}]
client_body_timeout 5s;
client_header_timeout 5s;
keepalive_timeout 15s;
send_timeout 5s;

limit_req_zone $binary_remote_addr zone=req_limit:10m rate=20r/s;
limit_req zone=req_limit burst=40 nodelay;
limit_conn_zone $binary_remote_addr zone=conn_limit:10m;
limit_conn conn_limit 20;
\end{lstlisting}

  \item \textbf{Prevención Activa de Intrusos con Suricata IPS (NFQUEUE):}
\begin{lstlisting}[style=terminal, caption={Reglas de Prevencion Activa en homelab.rules}, label={lst:reglas-ips}]
# Bloqueo inmediato de Inyeccion SQL (UNION SELECT)
drop tcp any any -> $HOME_NET $HTTP_PORTS (msg:"[BANK-IPS-DROP] Intento SQLi UNION SELECT"; flow:to_server,established; content:"UNION"; nocase; content:"SELECT"; nocase; classtype:web-application-attack; sid:1000001; rev:1;)

# Bloqueo inmediato de Path Traversal (/etc/passwd)
drop tcp any any -> $HOME_NET $HTTP_PORTS (msg:"[BANK-IPS-DROP] Path Traversal /etc/passwd"; flow:to_server,established; content:"etc/passwd"; nocase; classtype:web-application-attack; sid:1000003; rev:1;)
\end{lstlisting}

  \item \textbf{Efecto Tarpit y Descarte Silencioso (\textit{Silent Drop}):} La acción de descarte (\texttt{drop}) aplicada por Nftables y Suricata obliga a las herramientas del atacante (Nmap, bots) a esperar el agotamiento de sus temporizadores TCP, provocando los \texttt{TIMEOUT} y \texttt{ERROR: Script execution failed} documentados en la evidencia de escaneo.
\end{enumerate}

\section{Matriz Comparativa: Estado Inicial vs Estado Mitigado}

La \textbf{TABLA \ref{tab:comparativa-mitigacion}} resume el impacto de las soluciones técnicas implementadas, evidenciando la reducción drástica de la superficie de ataque y el fortalecimiento de la postura de seguridad global del sistema objetivo.

\begin{table}[H]
\centering
\caption{Comparativa Técnica: Estado Inicial (Vulnerable) vs Estado Mitigado}
\label{tab:comparativa-mitigacion}
\contenidoConFuente{%
  \begin{tablaAPA}
  \begin{tabular}{L{3.4cm}L{5.8cm}L{5.8cm}}
  \toprule
  \textbf{Aspecto de Seguridad} & \textbf{Estado Inicial (Vulnerable)} & \textbf{Estado Mitigado (Seguro)} \\
  \midrule
  \textbf{Cifrado SMTP (25, 587)} & Habilita intercambio anónimo Diffie-Hellman (\texttt{DH\_anon}) de 1024 bits. Vulnerable a MitM. & Suites anónimas (\texttt{aNULL}) prohibidas; forzado de cifrado TLS alto con DH de 2048 bits. \\
  \textbf{Base de Datos (5432)} & Expuesta a \texttt{0.0.0.0} en toda la LAN. & Confinada a red interna Docker; solo accesible por IPs administradoras. \\
  \textbf{Modo Debug Webmail} & Responde afirmativamente a verbos DEBUG (\texttt{http-aspnet-debug}). & Modo debug y volcado de errores desactivados en entorno de producción. \\
  \textbf{Cifrado Web en Tránsito} & Tráfico HTTP plano en puertos 8080, 8082, 3000. & Redirección obligatoria a HTTPS/TLS 1.3 con suites criptográficas seguras. \\
  \textbf{Cabeceras HTTP} & Ausencia de protecciones anti-Clickjacking, CSP y HSTS. & Cabeceras OWASP estrictas activas; divulgación de versión desactivada. \\
  \textbf{Escáneres y SQLi} & Respuesta normal a peticiones con inyecciones SQL. & WAF en Nginx bloquea inyecciones (403) e IPS Suricata descarta paquetes. \\
  \textbf{Escaneo Nmap (--script vuln)} & El atacante evalúa DoS y scripts agresivos sin restricción. & Efecto Tarpitting: descarte silencioso (\texttt{drop}) retiene al atacante en 99.58\% y causa timeouts. \\
  \bottomrule
  \end{tabular}
  \end{tablaAPA}%
}{Elaboración propia.}
\end{table}
